\subsection{INVERSE IMAGE OF AN EQUIVALENCE RELATION; INDUCED EQUIVALENCE RELATION}

Let \(\varphi\) be a mapping of a set \(E\) into a set \(F\) , and let \(S\) be an equivalence relation on \(F\) . If \(u\) is the canonical mapping of \(F\) onto \(F/S\) , the equivalence relation associated with the mapping \(u \circ \varphi\) of \(E\) into \(F/S\) is called the inverse image of \(S\) under \(\varphi\) ; if \(R\) is this relation, \(R\{x, y\}\) is equivalent to \(S\{\varphi(x), \varphi(y)\}\) , and the equivalence classes with respect to \(R\) are the inverse images under \(\varphi\) of the equivalence classes with respect to \(S\) which meet \(\varphi\langle E\rangle\) .

In particular, consider an equivalence relation R on a set E, and let A be a subset of E; then the inverse image of R under the injection j of A into E is called the equivalence relation induced by R on A, and is denoted by \(R_{A}\) .

The equivalence classes with respect to \(\mathbf{R}_{\mathbf{A}}\) are the traces on \(\mathbf{A}\) of the equivalence classes with respect to \(\mathbf{R}\) which meet \(\mathbf{A}\) . The injection \(j\) is obviously compatible with the relations \(\mathbf{R}_{\mathbf{A}}\) and \(\mathbf{R}\) ; the mapping \(h\) of \(\mathbf{A} / \mathbf{R}_{\mathbf{A}}\) into \(\mathbf{E} / \mathbf{R}\) induced by \(j\) on passing to the quotient with respect to \(\mathbf{R}_{\mathbf{A}}\) and \(\mathbf{R}\) is an injective mapping of \(\mathbf{A} / \mathbf{R}_{\mathbf{A}}\) into \(\mathbf{E} / \mathbf{R}\) : for if \(f\) (resp. \(g\) ) is the canonical mapping of \(\mathbf{E}\) onto \(\mathbf{E} / \mathbf{R}\) (resp. A onto \(\mathbf{A} / \mathbf{R}_{\mathbf{A}}\) ), then the relation \(h(g(x)) = h(g(x'))\) , where \(x \in \mathbf{A}\) and \(x' \in \mathbf{A}\) , is equivalent to \(f(x) = f(x')\) and therefore to \(g(x) = g(x')\) . The image \(h\langle \mathbf{A} / \mathbf{R}_{\mathbf{A}} \rangle\) is equal to \(f\langle \mathbf{A} \rangle\) . If \(k\) is the bijective mapping of \(\mathbf{A} / \mathbf{R}_{\mathbf{A}}\) onto \(f\langle \mathbf{A} \rangle\) which has the same graph as \(h\) , then \(k\) and its inverse are said to be canonical.

